% !TEX program = xelatex
% Main text: pages 1--4. Appendices: mechanism figure and evidence/QA notes.
\documentclass[11pt,a4paper,fontset=none]{ctexart}
\usepackage{hypothesis_generation}
\usepackage{amsmath,tabularx,pdflscape}
\geometry{left=19mm,right=19mm,top=22mm,bottom=20mm,headsep=7mm}
\setmainfont{Helvetica Neue}
\setsansfont{Helvetica Neue}
\setCJKmainfont{PingFang SC}
\setCJKsansfont{PingFang SC}
\setCJKmonofont{PingFang SC}
\definecolor{hypothesis1}{HTML}{0072B2}
\definecolor{hypothesis2}{HTML}{B46C00}
\definecolor{hypothesis3}{HTML}{8B4E7A}
\definecolor{evidencecolor}{HTML}{426D85}
\hypersetup{pdftitle={下一轮研究假设与实验设计},pdfauthor={本地研究设计记录},urlcolor=hypothesis1}
\fancyhead[L]{\small\color{hypothesis1}下一轮研究假设与实验设计}
\fancyhead[R]{\small S7 计划稿}
\fancyfoot[C]{\small 2026年9月6日\quad ·\quad \thepage}
\setlength{\parindent}{0pt}
\setlength{\parskip}{5pt}
\linespread{1.16}
\setlist{nosep,topsep=4pt,leftmargin=1.5em}
\tcbset{before skip=9pt,after skip=9pt}
\newcommand{\smallhead}[1]{\vspace{4pt}{\large\bfseries\color{hypothesis1}#1}\par}
\newcommand{\docref}[1]{{\small\color{darkgray}证据：#1。}}
\newcommand{\metric}{\ensuremath{Y(A,P,R;q)}}
\begin{document}

{\LARGE\bfseries\color{hypothesis1}下一轮研究假设与实验设计}\par
{\large 分开位置、关联路径和选图规则}\par
\textbf{状态：S7 尚未运行。}这是读取 S6 结果后提出的机制实验，不是 S6 的预注册预测，也不是新的独立测试集。

\begin{summarybox}[摘要：先查清参考照片为什么会改变]
同样20张历史照片，更新地图点的位置也会影响后续匹配。S6 的两种地图因此同时改变了坐标、来源关系和点数量，不能把结果只归因于“位置更准”。下一轮固定输入观测身份，重放关联事件，用四个地图和四种读出规则检查三条解释。

\textbf{H1 位置：}同一组观测仍归入同一点，只换最终坐标规则。\par
\textbf{H2 关联：}同一坐标规则，改换观测归组的历史；同时涉及点数、分区、出生属性和来源集合。\par
\textbf{H3 选择规则：}同一地图，改变候选限制或相近相机位置的抑制规则。

\textbf{三种解释可以共存；现在不宣布哪项主导。}
\end{summarybox}

\smallhead{已完成的证据与仍然缺少的证据}
\begin{tabularx}{\linewidth}{@{}p{33mm}X@{}}
\toprule
\textbf{S6 主测试} & \textbf{独立复算后的结果与范围}\\
\midrule
参考支持 & 91.907\% $\rightarrow$ 93.272\%，提高1.365个百分点；8张查询中4张换图。\\
简单对照 & 最近预测相机位置的4图为93.472\%；它没有NMS，尚非单因素解释。\\
几何与覆盖 & 共同像素MAE为367.526 $\rightarrow$ 347.882毫米；共同区域仅占有效目标约2.749\%。\\
反向结果 & 较稀疏采样下支持提高，MAE却从441.308升至447.767毫米。\\
\bottomrule
\end{tabularx}
\docref{S6结果、独立审查（附录B，证1--2）}

\textbf{审查已通过3453项。}其中2559项数值/值比较、432项完整性、461项元数据、1项源码顺序检查；检查次数不是样本数。内部查询状态原张量未保存，状态不变的核验仍以记录和代码为边界。

\smallhead{下一轮的固定边界}
沿用3块各24张照片：每块前20张历史、后4张只读查询。主测试仍为B1/B2的8张图，B0四张只作开发描述；主采样步长8，步长12作敏感性检查。不重新推理模型，不用实测深度或参考轨迹建图、选图。全部选择封存后才评分。

\textbf{当前强主张不成立：}未证明简单平均是新颖的优越方法，未证明支持代理提高等于视频改善，也没有完成完整VMem视频生成。下一步价值首先是可复核的机制诊断，仍需最近工作排重和独立场景验证。

\newpage
{\LARGE\bfseries\color{hypothesis1}H1｜最终位置是否足以改变决定？}\par
\begin{hypothesisbox1}[解释与成立前提]
在同一关联路径内，每个点包含完全相同的历史观测和来源照片；只把最终位置从首次观测坐标改为逐帧均值。位置变化可能改变投影格、遮挡先后和可见投票，进而改变所选4张照片。

\textbf{动机：}S6出现过换图，但同时改变了关联和点数；VMem原方法已用可见表面上的来源信息检索历史照片。两者都不能单独证明H1。\citep{vmem}

\textbf{前提：}固定路径内点数、出生法线/半径/颜色、来源集合和贡献帧次数必须严格相同；重放不得重新寻找邻居。
\end{hypothesisbox1}

\begin{predictionbox}[预测与可反驳条件]
\textbf{操作化预测：}固定关联路径A0或A1，切换P0/P1后，至少一个预先指定的主测试查询改变最终参考ID集合；记录变化如何从可见点进入投票与选择。

\textbf{反驳条件：}在两条固定路径、全部8张主测试查询中，最终集合均不变，则“只改最终位置足以改变本设置选图”的预测未成立。中间投票改变但最终集合不变，只说明选择规则吸收了变化。

\textbf{收益另判：}换图不保证支持上升。逐查询正、负、零变化全部报告；本轮没有预先假设可推广的提升幅度。
\end{predictionbox}

\smallhead{设计：两条固定路径内分别比较}
\begin{tabularx}{\linewidth}{@{}p{23mm}XX@{}}
\toprule
 & \textbf{P0：首次位置} & \textbf{P1：逐帧均值位置}\\
\midrule
A0 首写路径 & A0P0：必须复现S6首写 & A0P1：固定A0，只改位置\\
A1 均值路径 & A1P0：固定A1，只改位置 & A1P1：必须复现S6均值\\
\bottomrule
\end{tabularx}

每个贡献帧先形成一个质心，再按原递推规则等帧平均；不能改为全部像素等权平均。主要读出采用原选择器，其余三种读出用于检查解释是否依赖选择规则。以支持百分比$Y$记录每查询$Y(A_i,P_1)-Y(A_i,P_0)$，同时报告ID集合和顺序。

\begin{comparisonbox}[先验证，再解释]
A0P0与S6首写、A1P1与S6均值的全部数组、来源表和地图摘要必须严格一致；摘要不含不可重复的计时。重新录制的原Memory也须与对应重放一致。任何一项不一致即停止，不进入评分。
\end{comparisonbox}

\textbf{限制：}这是固定关联下的反事实重放，不等于新的在线建图算法。共同像素误差不能代表完整表面；即使H1成立，也不排除H2或H3。未出现变化只限制本批数据，不能证明位置永远无关。\par
\docref{S7协议；S6结果和独立审查（附录B，证1--3）}

\newpage
{\LARGE\bfseries\color{hypothesis2}H2｜历史观测怎样归组是否关键？}\par
\begin{hypothesisbox2}[解释与成立前提]
在线更新位置会改变后来观测匹配到哪个点。即使最后采用同一种位置规则，两条关联历史也可能产生不同的点组和来源集合，从而改变哪些历史照片获得投票。

\textbf{这里的“关联路径效应”包含：}点数量、观测分区、点出生时的法线/半径/颜色等属性，以及来源照片集合。\textbf{不能把它缩写成只改来源ID，也不预先断言关联更正确。}

\textbf{动机：}S6两臂的点数确实不同，且支持与几何误差并不总同向；这提供检验动机，没有证明关联主导。
\end{hypothesisbox2}

\begin{predictionbox}[预测与可反驳条件]
\textbf{操作化预测：}固定P0或P1，只替换A0/A1后，至少一个主测试查询的参考ID集合改变。记录首次导致可见来源或候选变化的执行步骤。

\textbf{反驳条件：}两种固定位置规则下，全部8张主测试查询的最终集合均不变，则该干预未显示H2的选图作用。若位置干预已产生变化，也不能只凭S6对角两臂宣布“关联是唯一原因”。

\textbf{不设“主导”胜负：}本轮不以少量查询比较绝对效应后选唯一解释；ID变化和支持差值分别保留。
\end{predictionbox}

\smallhead{设计：先记住每个观测的固定身份}
每个接受的历史观测用“帧编号、原224裁剪图像素$u,v$”标识。分别按首写路径A0和在线均值路径A1记录匹配目标；未匹配观测按原次序出生。两个路径接受的观测身份与原几何顺序必须一致。同帧新出生点不可被本帧后续观测匹配。

随后只重放既有归属，得到四臂；固定P0比较A1P0与A0P0，固定P1比较A1P1与A0P1。跨路径点数可能不同，\textbf{不按最终点编号硬配对，也不直接交换来源表。}

\begin{comparisonbox}[交互只作描述，不能冒充唯一因果归因]
每个查询的支持交互量为
\[
  (Y_{A1P1}-Y_{A1P0})-(Y_{A0P1}-Y_{A0P0}).
\]
它描述两种干预是否相互依赖。零交互不证明只有一个机制；非零交互也不证明哪条路径在自然场景中主导。
\end{comparisonbox}

\textbf{限制：}H2仍同时改变多种属性。若要单独研究来源标签，需在同一地图身份上另立干预；当前四臂不提供该结论。单房间的有限查询不足以支持一般化新方法，所有下降、零变化和开发集结果均保留。\par
\docref{S6点数表；S7协议；候选贡献反证审查（附录B，证1、3--4）}

\newpage
{\LARGE\bfseries\color{hypothesis3}H3｜选择规则是否放大或吸收变化？}\par
\begin{hypothesisbox3}[解释与成立前提]
同一地图提供的候选和票数，经过排序、去除相近相机位置的NMS后，未必选出同一组参考。H3检查候选限制与NMS怎样改变最终决定；它可以与位置或关联效应同时成立。

\textbf{动机：}S6最近预测姿态4图的支持稍高，但与原选择器还存在NMS等差异；不能据此单独责怪某个模块。VMem已有来源索引及NMS，AnchorWeave也已有覆盖驱动检索；“换一种选图规则”本身不是新颖性证明。\citep{vmem,anchor}

\textbf{前提：}比较时地图、查询相机、参考数量4、距离精度和并列顺序固定；只改变表中明确指定的因素。
\end{hypothesisbox3}

\smallhead{设计：同一地图的四种读出}
\begin{tabularx}{\linewidth}{@{}p{33mm}XX@{}}
\toprule
\textbf{读出} & \textbf{候选} & \textbf{最后选择}\\
\midrule
official & 原展开候选，多重计数保留 & 原NMS和放宽规则\\
candidate\_no\_nms & 与official完全相同 & 排序后前4个唯一ID\\
all20\_nms & 20张历史各一次 & 同样NMS和放宽规则\\
all20\_no\_nms & 20张历史各一次 & 排序后前4个唯一ID\\
\bottomrule
\end{tabularx}

\begin{predictionbox}[预测与可反驳条件]
\textbf{操作化预测：}同地图下，至少一个主测试查询因移除NMS或放宽候选而改变最终集合；保存候选排序、每轮拒绝原因及阈值放宽/回填记录，以定位变化。

\textbf{反驳条件：}所有预先指定同地图对照均保持集合不变，则本设置未显示H3的最终选图作用。若换图但支持下降，则否定“该改动改善本代理”的方向性判断，不能只展示有利查询。
\end{predictionbox}

\smallhead{实施前必须排除两个伪消融}
\textbf{第一，不能直接关闭官方NMS开关。}该代码分支还会强制加入最后一张历史并改变阈值。辅助实现须保留候选展开顺序、float32距离排序和并列语义；查询先经过官方单姿态四元数规范化。带NMS路径逐次复现官方最终顺序后才能使用。

\textbf{第二，不能悄悄补足参考图。}任一路径不能得到4个不同的0--19历史ID时立即停止，报告设置局限。全20候选的两种读出不依赖地图，应在四地图间严格一致。

\textbf{限制：}支持是50毫米阈值下的实测几何可见性代理，不含视频美观或生成一致性。四种读出、两种采样密度、多个像素不是独立样本。只报逐查询变化与主8图均值，不做总体显著性检验。\par
\docref{S7协议；S6简单对照；VMem方法（附录B）}

\clearpage
\begin{landscape}
{\Large\bfseries\color{hypothesis1}附录A｜可编辑机制图：计划消融，不是新结果}\par
\vspace{3mm}
\begin{center}
\includegraphics[width=\linewidth,height=0.81\textheight,keepaspectratio]{下一轮机制图.pdf}
\end{center}
\small 此图由本机Draw.io实际导出。蓝、橙、紫与H1/H2/H3标签同时区分解释。外部实测深度和GT轨迹只连接评分；H2也会改变点数量。完整文字可在同目录原生\texttt{.drawio}文件中编辑。
\end{landscape}

\clearpage
{\LARGE\bfseries\color{hypothesis1}附录B｜执行约束、证据与制作说明}\par
\smallhead{共同执行约束与研究判断}
\begin{itemize}
\item 计划共6案例$\times$4地图$\times$4查询$\times$4读出，即384个选择条件；仍只有12张不同查询、主测试8张。检索宽度只用160；S6的160/320完全相同，不当作重复。
\item 所有选择与渲染数据先封存SHA，再读已核验的S6测量数组。使用前逐值对照S6与独立审查的支持/目标并核源哈希；明确属于复用已验证测量，不称重新从PNG计算。
\item 几何评分沿用S6尺度与GT约定，重算112网格最近正深度。四臂共同像素、各自覆盖及对角两臂的共同像素均报告；空共同区域记为空，不删查询。
\item 保留原投票器按可见像素累加、每个来源首项双加的执行语义。位置平均的“每帧一个质心”不等于把检索投票改成每帧一票。
\item 先做人工小例核程序语义，再做真实观测重放；人工例不计研究样本。结果须由另一实现从观测ID/事件重建地图、来源与全部支持；不得导入本轮实现自证。
\end{itemize}
\textbf{质量判断：}三假设均可在本机固定数组上检验，且有设置内的反驳条件。H1控制更直接；H2仍包含多属性联动；H3最需防止代码开关带入额外变化。三者的新颖性与跨场景解释力都尚未确立。

\smallhead{项目证据：优先于早期文稿中的状态用语}
{\small
\begin{description}[style=nextline,leftmargin=11mm,itemsep=2pt]
\item[证1] \texttt{docs/S6\_RESULTS.md}：已完成的组件实验与全部限制。
\item[证2] \texttt{docs/S6\_INDEPENDENT\_AUDIT.md}及\texttt{results/S6\_independent\_audit/verification.json}：3453项通过，区分真实复算、元数据和状态记录边界。
\item[证3] \texttt{docs/S7\_EVENT\_REPLAY\_PROTOCOL.md}：本报告所述四臂、读出、封存与停机条件的详细规范。
\item[证4] \texttt{docs/NOVELTY\_CANDIDATE\_REVIEW.md}：候选A为本轮机制诊断；候选B效率和候选C冲突关系尚未实现，不混入本轮H1--H3。
\end{description}}

\smallhead{两项直接相关的已核验原始文献}
\vspace{-3mm}
\begingroup\small
\renewcommand{\section}[2]{}%
\begin{thebibliography}{2}
\bibitem[Li et al.(2025)]{vmem} Runjia Li, Philip Torr, Andrea Vedaldi, Tomas Jakab. \textit{VMem: Consistent Interactive Video Scene Generation with Surfel-Indexed View Memory}. arXiv:2506.18903v3, 2025. \href{https://arxiv.org/html/2506.18903v3}{原文：方法3.1及附录A}。来源索引、可见投票和NMS已有；本轮研究其执行机制，不声称发明这些组件。
\bibitem[Wang et al.(2026)]{anchor} Zun Wang, Han Lin, Jaehong Yoon, Jaemin Cho, Yue Zhang, Mohit Bansal. \textit{AnchorWeave: World-Consistent Video Generation with Retrieved Local Spatial Memories}. arXiv:2602.14941v1, 2026. \href{https://arxiv.org/html/2602.14941v1}{原文：方法3.2--3.4}。已有局部几何记忆、覆盖检索和多锚点条件；不能仅凭这些概念声称创新。
\end{thebibliography}
\endgroup

{\small\textbf{制作说明：}采用本地Scientific Hypothesis Generation技能的竞争解释、可反驳预测与配对设计方法，按本轮任务保留四页主文和必要附录，没有为50条引用的模板目标补造文献。使用其本地样式并调整中文字体与颜色；图实际改用可编辑Draw.io，\textbf{未运行NanoBanana/Gemini或Claude模型流程}。本报告由本机XeLaTeX编译并逐页渲染检查，命令、返回码、文件哈希与QA随源文件交付。S7尚未运行；源协议与旧结果均不因排版修改。}
\end{document}
